%% kapa: exam=1 round=21 year=2010 period=2 form=A subject=accounting subject_ko=회계학 seq=31 number=31
%% source: Q-Net 감정평가사 기출문제 (cst00309) artlSeq=5024181 file=extracted/1008/5024181/2010년 제21회 감정평가사 1차시험 기출문제/2교시.hwp
%% conversion: kordoc 4.24.51 -> tools/gen1.py (auto); review=proofread-gpt-sol; flags=none
\begin{question}{31}
프린터를 생산ㆍ판매하는 (주)한국의 최대생산능력은 연 12,000대이고, 정규시장에서 연간 판매량은 10,000대이다. 단위당 판매가격은 ￦100,000이고, 단위당 변동제조원가는 ￦60,000이며, 단위당 변동판매비와 관리비는 ￦10,000이다. (주)한국은 (주)서울로부터 프린터 4,000대를 단위당 ￦70,000의 가격으로 구입하겠다는 1회성 특별주문을 받았다. (주)한국은 올해 생산능력을 변경할 계획이 없다. (주)한국의 판매비와 관리비는 모두 변동비인데, (주)서울의 주문을 받아들이는 경우 이 주문과 관련된 (주)한국의 판매비와 관리비는 75\%가 감소할 것으로 추정된다. (주)한국이 동 주문을 수락하기 위하여 기존시장의 판매를 일부 포기하기로 한다면 증분손익은 얼마인가? (단, 기초ㆍ기말재고는 없음)
\begin{choices}
  \item ￦7,500,000 감소
  \item ￦30,000,000 감소
  \item ￦7,500,000 증가
  \item ￦30,000,000 증가
  \item ￦0
\end{choices}
\end{question}
\answer{2}
